\documentclass[11pt]{ctexart}
\usepackage[a4paper,top=2cm,bottom=2cm,left=2.2cm,right=2.2cm]{geometry}
\usepackage{fancyvrb}
\usepackage{xcolor}
\usepackage{colortbl}
\usepackage{fontspec}
\usepackage[most]{tcolorbox}
\usepackage{booktabs}
\usepackage{enumitem}
\usepackage{titlesec}
\usepackage{hyperref}

\setmonofont{Consolas}
\setCJKmonofont{SimHei}
\setlength{\parindent}{0pt}
\setlist{nosep,leftmargin=1.4em}

% ── 嵌入式（一）主题色：深琥珀棕（PCB/铜色系，视觉五册之外的新族） ──
\definecolor{deep}{RGB}{92,58,20}        % 主色：深琥珀棕
\definecolor{accent}{RGB}{162,104,36}    % 强调：琥珀
\definecolor{keyline}{RGB}{90,130,200}   % 知识点：蓝
\definecolor{keybg}{RGB}{238,244,255}
\definecolor{warnline}{RGB}{200,120,60}  % 警告：橙
\definecolor{warnbg}{RGB}{255,243,235}
\definecolor{okline}{RGB}{26,116,88}     % 成果：绿
\definecolor{okbg}{RGB}{240,252,244}
\definecolor{tipline}{RGB}{170,140,92}   % 提示：浅琥珀
\definecolor{tipbg}{RGB}{250,246,238}
\definecolor{codebg}{RGB}{249,249,252}   % 代码底
\definecolor{rowtint}{RGB}{248,242,230}  % 表头浅琥珀

\titleformat{\section}{\large\bfseries\color{deep}}{\thesection}{0.6em}{}
\titleformat{\subsection}{\normalsize\bfseries\color{accent}}{\thesubsection}{0.5em}{}
\titleformat{\part}[display]{\centering\Huge\bfseries\color{deep}}{}{0pt}{\Huge}
\titlespacing{\section}{0pt}{14pt}{6pt}
\titlespacing{\subsection}{0pt}{10pt}{4pt}

\newtcolorbox{keybox}[1]{enhanced, colback=keybg, colframe=keyline,
  boxrule=0.6pt, arc=1.5mm, left=3mm, right=3mm, top=2mm, bottom=2mm,
  title={\bfseries #1}, fonttitle=\small\bfseries, coltitle=white,
  colbacktitle=keyline}
\newtcolorbox{warnbox}[1]{enhanced, colback=warnbg, colframe=warnline,
  boxrule=0.6pt, arc=1.5mm, left=3mm, right=3mm, top=2mm, bottom=2mm,
  title={\bfseries #1}, fonttitle=\small\bfseries, coltitle=white,
  colbacktitle=warnline}
\newtcolorbox{okbox}[1]{enhanced, colback=okbg, colframe=okline,
  boxrule=0.6pt, arc=1.5mm, left=3mm, right=3mm, top=2mm, bottom=2mm,
  title={\bfseries #1}, fonttitle=\small\bfseries, coltitle=white,
  colbacktitle=okline}
\newtcolorbox{tipbox}[1]{enhanced, colback=tipbg, colframe=tipline,
  boxrule=0.5pt, arc=1.5mm, left=3mm, right=3mm, top=1.5mm, bottom=1.5mm,
  title={\bfseries #1}, fonttitle=\small\bfseries, coltitle=white,
  colbacktitle=tipline}

\newcommand{\code}[1]{\texttt{#1}}
\DefineVerbatimEnvironment{Verilog}{Verbatim}{frame=single,fontsize=\small}

\begin{document}

% ==================== 封面 ====================
\begin{titlepage}
\centering
\vspace*{2.2cm}
{\color{deep}\Huge\bfseries 全链集成与仿真设施}\\[6pt]
{\color{deep}\Large\bfseries 嵌入式对接教程（一）}\\[14pt]
{\large 2026 FPGA 创新设计赛道（安路）赛题二}\\[6pt]
{\large 嵌入式方向 · 04\_embedded}\\[20pt]

\begin{tcolorbox}[enhanced, colback=codebg, colframe=accent, boxrule=0.8pt,
  arc=2mm, width=12.5cm, left=4mm, right=4mm, top=3mm, bottom=3mm]
\centering
{\small 读者：嵌入式方向队员（零基础，建议先浏览视觉 Step1 教程）\\
目标：看懂官方 ISP 链路，把视觉模块装进去，并把整链仿真跑起来\\
前置要求：会 Verilog、跑过 ModelSim\\
预计学习时间：8～12 天（每天 2～3 小时）}
\end{tcolorbox}

\vspace{8mm}
\begin{tcolorbox}[enhanced, colback=tipbg, colframe=tipline, boxrule=0.5pt,
  arc=2mm, width=12.5cm, left=4mm, right=4mm, top=2mm, bottom=2mm]
\centering
{\small\color{deep} \textbf{这一册解决什么}\\
视觉方向交付的是"算法模块"；嵌入式方向负责把它们\textbf{装进官方}
\textbf{链路}、并在没有板子的条件下\textbf{证明装对了}——本册讲集成与仿真。\\
外设驱动、摄像头配置与固化见第二册。}
\end{tcolorbox}

\vfill
{\small\color{black!50} fpga-ptz-tracker 教程系列 · 2026-09}
\end{titlepage}

\tableofcontents
\newpage

% ============================================================
\part{从零学习教程}

% ============================================================
\section{你在造什么}

\subsection{一句话}

\textbf{把视觉模块装进官方摄像头链路，并在 ModelSim 里把"从 RAW 像素}
\textbf{到目标坐标"的整条路跑通——不用板子。}

\subsection{对接地图}

\begin{Verilog}
虚拟/真实 RAW10 传感器流 (40bit/拍)
      |
      v
+---------------------+
| isp_top（官方 ISP）  |  RAW10 解包 -> BLC -> demosaic -> AWB
| 明码 RTL             |  (MIPI PHY 与 DDR/HDMI 是加密黑盒, 不含)
+---------------------+
      |  128bit RGB
      v
+---------------------+
| data128_96（官方）   |  纯组合重排: 128bit -> 96bit (4像素/拍)
+---------------------+
      |  96bit AXI4-Stream  <- 视觉与嵌入式之间的"合同"
      v
+---------------------+
| vision_top（视觉）   |  检测 / OSD / 云台 / 声光
+---------------------+
      |
      +--> OSD 流 (帧转储)   +--> PWM x2 (云台)
\end{Verilog}

\subsection{做完是什么效果}

仿真的最后一帧画面被转储成图片（\code{sim/sys\_osd\_0.hex}），
红块被检测链套住、OSD 文字齐全；日志打出
\code{TB\_RESULT: PASS}——整链对接在无硬件条件下具备可复现的证据。

\begin{okbox}{学完本册你能做到}
\begin{itemize}
  \item 讲清官方 ISP 链路每级的职责与数据位宽
  \item 解释为什么视觉模块挂在 AWB 之后的 96bit 点上
  \item 为"仿真跑不了的门级网表/TD 原语"写行为替身，并守住替身纪律
  \item 读懂并运行 system\_top + tb\_system\_integration 整链仿真
  \item 说出三个"仿真与上板必然不同"的点，上板时先查它们
\end{itemize}
\end{okbox}

% ============================================================
\section{阶段 0：认识官方 ISP 链路（2 天）}

\subsection{一条链，两级分工}

官方例程（Lab1 通路）把摄像头数据一路加工到屏幕：

\begin{Verilog}
MIPI 差分信号 -> [PHY 硬核/加密] -> 2-lane 字节流
   -> raw10 解包 (拆成每像素 10bit RAW)
   -> BLC (黑电平校正)
   -> demosaic (Bayer 去马赛克: 每像素只有 R/G/B 一个分量 -> 三通道)
   -> AWB (自动白平衡)
   -> 128bit RGB 流 (8 像素/拍)
\end{Verilog}

我们能仿真的是\textbf{明码部分}（raw10 解包之后到 AWB 输出）；
上游 PHY 硬核与下游 DDR/HDMI 是\textbf{黑盒}（加密或依赖物理接口）。

\begin{keybox}{"明码/黑盒"的工程含义}
明码 = 能在 ModelSim 里编译仿真、能读代码定位问题；
黑盒 = 只能上板验证。本册所有仿真都建立在"明码段可信"之上，
黑盒段的验证任务明确交接给上板阶段（见检查点清单）。
\end{keybox}

\subsection{RAW10 与 Bayer 阵列}

传感器每像素输出 10bit 的\textbf{单色}强度，按 BGGR 图案排布：

\begin{Verilog}
偶行:  B G B G B G ...
奇行:  G R G R G R ...
(2x2 重复单元; demosaic 的任务就是"猜"出每个像素缺的两个分量)
\end{Verilog}

40bit 一拍正好打包 4 个像素（4 x 10bit），行的长度 1280/4 = 320 拍。

\subsection{练习}

\begin{enumerate}
  \item 手算：1280x720 的一帧，RAW10 输入共有多少拍？
  \item 思考：为什么红块的"红"在 RAW 里只有 R 位点的像素偏红？
    （提示：先看 BGGR 排布）
  \item 找出官方例程里 MIPI PHY 文件夹，确认它是加密的（\code{.enc}）
\end{enumerate}

% ============================================================
\section{阶段 1：96bit 契约与挂载点（1～2 天）}

\subsection{契约三要素（+握手）}

视觉模块消费的是 96bit AXI4-Stream，规则在
\code{06\_interfaces/current/IF-MIPI.md} 里钉死：

\begin{Verilog}
tdata[95:0]   一拍 4 像素 RGB888 (每像素 24bit)
tuser         帧首拍 = 1 (sof)
tlast         行尾拍 = 1 (eol)
tvalid/tready 默认不反压 (tready 恒 1)
\end{Verilog}

任何"塞进这条流"的模块都必须遵守这三要素；视觉的 detect\_top
甚至把它当\textbf{纯旁路}处理（输入原样输出，检测在侧支进行）。

\subsection{为什么挂在 AWB 之后}

\begin{keybox}{挂载点的三个理由}
\textbf{1. 颜色已校正}：AWB 拉平了光源色偏，红色球在白炽灯下也能稳定
落在红窗里——挂在 AWB 之前，阈值就得跟着灯光漂。\\
\textbf{2. 位宽友好}：AWB 输出就是 128bit RGB，经 data128\_96 重排成
96bit（4 像素/拍）正好是视觉链路的舒适粒度。\\
\textbf{3. 零侵入}：这一级是纯数据通路，旁路挂载不干扰官方时序。
\end{keybox}

\subsection{data128\_96：官方模块原样复用}

128bit（8 像素/拍）到 96bit（4 像素/拍）只是\textbf{纯组合重排}，
官方模块一行不改直接用。这是"契约层"的证据：视觉侧 14 项回归
证明它完全兼容这条流。

\subsection{练习}

\begin{enumerate}
  \item 画出 128bit 拍拆成两拍 96bit 的对应关系（哪个像素去哪一拍）
  \item 思考：如果 AWB 关掉（锁死增益），检测阈值会如何变化？
\end{enumerate}

% ============================================================
\section{阶段 2：仿真替代品哲学（2～3 天）}

\subsection{问题：有些文件仿真器根本编不过}

官方工程里有两类"仿真毒药"：

\begin{itemize}
  \item \textbf{门级网表}：\code{divider\_gate.v} 是 AL\_* 库单元拼出来的
    7355 行网表，ModelSim 里没有这套库；
  \item \textbf{厂商原语}：官方 blk\_mem 仿真 RTL 例化 TD 的
    \code{PH1P\_LOGIC\_ERAM} 原语，ModelSim 也不认识。
\end{itemize}

直接把官方全套塞给 ModelSim？第一次尝试炸了 \textbf{58840 个错误}。

\subsection{解法：行为级替身（sim-only stand-in）}

为每个"编不过的"文件写一个\textbf{行为等价}的替身：仿真用替身、
TD 综合用真文件——两条腿走路，互不干扰。

\begin{Verilog}
divider_gate.v (AL_* 门级网表)  --仿真-->  divider_gate_sim.v (行为: 16 拍除法)
                                --综合-->  divider_gate.v (原样)

PH1P_LOGIC_ERAM (TD 原语)       --仿真-->  ph1p_logic_eram_stub.v (双口 RAM)
                                --综合-->  真原语
\end{Verilog}

\begin{warnbox}{替身三纪律（写进 filelist 才算数）}
\textbf{1. 端口/参数逐一对照}：替身的接口清单必须从"官方实例化现场"
抄下来（含 ECC 端口的\textbf{方向}：\code{sbiterr/dbiterr} 是输出、
\code{*inj} 是输入——抄反了编译能过、行为全错）。\\
\textbf{2. 行为对齐语义}：替身的延迟、除零行为、读写模式要和真件一致
（不要求逐拍一致，但"能跑通流程"必须一致）。\\
\textbf{3. 严禁进 TD 工程}：\code{filelist.f} 单列一节
"仿真专用（禁止加入 TD 工程）"，上板综合时一旦带进去，
同模块双定义或行为替换，后果随机。
\end{warnbox}

\subsection{练习}

\begin{enumerate}
  \item 打开 \code{divider\_gate\_sim.v}，回答：start 后多少拍出 done？
    除零时商是多少？
  \item 找一个官方 blk\_mem 的实例化行，数出它用了多少个端口，
    和替身模块的端口列表对一遍
  \item 思考：替身的延迟和真件不同，为什么整链仿真仍然可信？
    （提示：功能级 vs 像素级）
\end{enumerate}

% ============================================================
\section{阶段 3：system\_top 集成（2 天）}

\subsection{顶层结构}

\code{system\_top.v} 把一切装进一个盒子：

\begin{Verilog}
I_raw10_* --> isp_top --> 128bit --> data128_96 --> 96bit
                                                  |
                    +-----------------------------+
                    v
              vision_top (检测/OSD/云台/声光)
                    |
      +-------------+--------------+-------------+
      v             v              v             v
   O_osd_*     O_pwm_x/y      O_led/buzzer   uart_telemetry
   (帧转储)     (舵机)          (声光)         (遥测)
\end{Verilog}

按键进消抖模块，输出两个\textbf{触发脉冲}：KEY0 切阈值组、
KEY1 翻转跟踪使能——十行 glue 逻辑，不写更复杂的状态机。

\subsection{两种"接线"姿势：端口 vs 层次化引用}

\begin{itemize}
  \item \textbf{端口连接}（正规）：所有跨模块信号走端口；
  \item \textbf{层次化引用}（例外）：TB 里读 DUT 内部信号用
    \code{u\_sys.u\_vision.t\_cx} 直接穿透——仿真专用，方便；
    但 RTL 内部互相引用要克制。
\end{itemize}

本工程里有一处 RTL 内部也用了层次化引用（遥测要拿 vision 的内部坐标），
属于"顶层特权"：只在最顶层允许，子模块之间严禁。

\subsection{两个真实踩坑}

\begin{warnbox}{坑 1：端口方向看名字想当然}
\code{isp\_top} 的 \code{I\_tready} 听起来是"输入"，实际是
\textbf{输出}（对上游的就绪指示）。接成常量 0 直接报错。
教训：\textbf{接黑盒/明码模块的端口前，先看模块头部的方向声明}，
名字前缀（I\_/O\_）不等于方向，尤其官方代码。
\end{warnbox}

\begin{warnbox}{坑 2：先赋值后声明}
system\_top 初版里 vision 内部信号"先赋值后声明"——Verilog 隐式
wire 规则会把它们悄悄建成 1bit 新信号，编译报重复声明。
修复：\textbf{声明前置}，统一集中在模块头部。
（和 Step5 的 m\_* 顺序坑同根：Verilog 隐式声明是事故之源。）
\end{warnbox}

\subsection{练习}

\begin{enumerate}
  \item 数一数 system\_top 里例化了几个模块，各属于官方/视觉/嵌入式哪一类
  \item 手写 KEY0 切阈值组的 ten-line glue（3bit 计数器 + 触发脉冲）
  \item 思考：为什么遥测用层次化引用而不是在 vision\_top 上加一个端口？
\end{enumerate}

% ============================================================
\section{阶段 4：跑通整链对接仿真（2 天）}

\subsection{虚拟 RAW10 传感器}

没有摄像头，就用 Python 生成一个"假摄像头"：BGGR 灰底 + 红块
（R 位点 900、G/B 位点 64）+ 0.2\% 确定性噪声，
按官方输入格式打包成 40bit/拍（\code{raw10\_beats.hex}），
标志序列单独一行一个数（\code{raw10\_meta.txt}，bit1=sof、bit0=eol）。

\begin{keybox}{生成器的价值在于"可复现"}
噪声用固定种子（\code{default\_rng(500)}），每次生成的文件逐字节
相同——仿真失败可以重放，金标准才站得住。
（生成器本身也曾有 bug：heredoc 转义吃掉替换，只出一帧多——
改行级补丁后修复，见 Part II 复盘。）
\end{keybox}

\subsection{tb\_system\_integration 的四条判据}

\begin{center}
\begin{tabular}{lll}
\toprule
\rowcolor{rowtint}\textbf{判据} & \textbf{内容} & \textbf{性质} \\
\midrule
A & ISP 几何：每帧 tuser 恰 1 次、tlast 恰 720 次 & 结构正确 \\
B & 检测锁定：valid=1 且坐标在红块 ±40px 内 & 功能正确 \\
C & 云台方向：块在左半屏则 duty\_x < 1500 & 方向正确 \\
D & UART 遥测：收到完整 A5..5A 帧且坐标一致 & 外设正确 \\
\bottomrule
\end{tabular}
\end{center}

注意判据 B 是\textbf{功能级}（±40px 容差）不是像素级——
因为 AWB 会自适应改变颜色，检测结果本就不是逐位可复现的。
验收口径要跟着链路性质走，这是"分层证明"的一个实例。

\subsection{重要发现：tuser 链在纯 RTL 仿真不激活}

第一次跑整链，官方 demosaic 的 tuser（帧同步标志）链\textbf{全灭}：
bilinear 5 拍延迟线输入无脉冲（tlast 正常数到 2424 次）。
分析结论：官方 demosaic 的 tuser 生成依赖\textbf{真实 CSI 时序或
TD 仿真库环境}，纯 ModelSim 行为级复现不出来。

\begin{warnbox}{这不算失败，算"界碑"}
工程处置：把全链行为仿真\textbf{明确列为 TD 阶段验证项}，
并把这条发现写进上板检查点（"上板有数据无帧同步先查 tuser 链"）。
无硬件条件下，该做的不是硬凑一个假通过，而是\textbf{把边界如实画出来}。
\end{warnbox}

\subsection{练习}

\begin{enumerate}
  \item 算一算：4 帧 RAW10 激励文件共多少拍？多少 MB（hex 文本）？
  \item 手推判据 A 的检查逻辑：tuser 计数在"帧边界"怎么重置？
  \item 思考：为什么判据 B 的容差是 40px 而不是 5px？
\end{enumerate}

% ============================================================
\section{阶段 5：把结论写成可交接的材料}

\subsection{两份 README}

\begin{itemize}
  \item \code{README\_嵌入式方向报告.md}：四步交付总览 + \textbf{分层证明}
    （契约层/挂载层/行为层/外设层各自证到了什么）+ 发现 + 遗留；
  \item \code{README\_分工与移交.md}：给嵌入式队友的挂载清单
    （挂哪几个模块、加哪些引脚、联合验收勾选表）。
\end{itemize}

\begin{keybox}{分层证明的写法（答辩直接可用）}
不要说"对接完成了"。说：\textbf{契约层已证}（视觉 14 项回归证明兼容
96bit 契约）、\textbf{挂载层已证}（elaboration 零错误）、
\textbf{行为层部分}（整链仿真受 demosaic tuser 限制，列 TD 阶段）、
\textbf{外设层已证}（UART 逐字节验证）。每一层都有证据文件和判据——
这是工程可信度的写法。
\end{keybox}

\subsection{练习}

\begin{enumerate}
  \item 用一句话给每层证明写出"主张 + 证据文件"
  \item 列出你认为还有哪个层没被覆盖（提示：黑盒段）
\end{enumerate}

% ============================================================
% 嵌入式对接教程（一）· Part II：工程实现详解
% 由 emb1_tutorial.tex \input 引入（共用同一份导言区与主题色）
% ============================================================

\part{工程实现详解}

\begin{tipbox}{这一部分写什么}
把集成与仿真设施逐文件讲透：官方链路的接入细节、两个行为替身
（divider\_gate\_sim 52 行 / ph1p\_logic\_eram\_stub 160 行）、
system\_top 顶层、虚拟 RAW10 传感器与整链 TB，
以及 8 轮迭代里属于本册的 5 轮复盘。
\end{tipbox}

% ============================================================
\section{工程总览}

\subsection{本册涉及的文件}

\begin{center}
\begin{tabular}{llr l}
\toprule
\rowcolor{rowtint}\textbf{文件} & \textbf{类型} & \textbf{行数} & \textbf{职责} \\
\midrule
\code{rtl/system\_top.v} & RTL & 154 & 集成顶层（仿真/TD 双用） \\
\code{rtl/divider\_gate\_sim.v} & 替身 & 52 & 门级除法器的行为级替身 \\
\code{rtl/ph1p\_logic\_eram\_stub.v} & 替身 & 160 & TD ERAM 原语行为替身 \\
\code{tb/tb\_system\_integration.v} & 测试 & 163 & 整链对接四判据 \\
\code{sim/isp\_top.v} 等 & 官方副本 & --- & 官方明码 RTL（仿真用） \\
\code{sim/raw10\_beats.hex} & 向量 & --- & 4 帧虚拟传感器（230400 拍/帧） \\
\midrule
\code{tools/gen\_raw10\_sensor.py} & 脚本 & --- & 虚拟传感器生成器 \\
\bottomrule
\end{tabular}
\end{center}

\subsection{四步交付与分层证明}

嵌入式交付的四步（环境/驱动/配置/联调）在报告里各占一格；
本册覆盖第（1）步的仿真环境与第（4）步的集成联调，
第二册覆盖第（2）（3）步（驱动与摄像头配置）。分层证明的对应：

\begin{center}
\begin{tabular}{lll}
\toprule
\rowcolor{rowtint}\textbf{层} & \textbf{主张} & \textbf{证据} \\
\midrule
契约层 & 视觉兼容 96bit 契约 & tb\_vision\_top 14 项回归 \\
挂载层 & 官方 ISP + 视觉可共存于一个顶层 & elaboration 零错误 \\
行为层 & 整链数据流贯通（受 demosaic 限制） & tb\_system\_integration 判据 A/C \\
外设层 & UART 遥测逐字节正确 & tb\_uart\_telemetry（第二册详述） \\
\bottomrule
\end{tabular}
\end{center}

% ============================================================
\section{官方链路的接入细节}

\subsection{isp\_top 实例化的两个坑}

\begin{Verilog}
isp_top u_isp (
    ...
    .I_tvalid  (I_raw10_valid & isp_tready),  // 与上游就绪相与
    .O_tready  (isp_tready),                  // 它是输出: 对上游指示
    .I_tready  (isp_up_tready)                // 也是输出(命名反向)
);
\end{Verilog}

坑一：\code{I\_tready} 实际方向是输出，接常量直接编译报错；
坑二：\code{I\_tvalid} 要与 \code{isp\_tready} 相与——真实流控语义，
上游数据只有在 ISP 就绪时才有效。\textbf{读官方代码先读端口注释，
不要相信前缀命名}。

\subsection{RAW10 进、128bit 出}

数据传输的实际形状：

\begin{Verilog}
40bit/拍 (4 像素 x 10bit) x 320 拍/行 x 720 行 = 230400 拍/帧
   -> isp_top 内部: 解包/黑电平/去马赛克/白平衡 (整拍流水)
   -> 128bit/拍 (8 像素 x 24bit RGB) 输出
   -> data128_96 纯组合重排 -> 96bit (4 像素)
\end{Verilog}

% ============================================================
\section{仿真替身详解}

\subsection{divider\_gate\_sim：一个"假装很慢"的除法器}

真件是 7355 行 AL\_* 门级网表；替身用 FPGA 也不擅长、但仿真器
支持的行为级除法：

\begin{Verilog}
if (!busy && start && !done) begin
    num_q <= numerator; den_q <= denominator; busy <= 1; wait_cnt <= 0;
end else if (busy) begin
    if (wait_cnt == 6'd15) begin            // 16 拍"伪装"门级延迟
        if (den_q != 0) begin
            quotient  <= num_q / den_q;      // 行为级除法(仿真器算)
            remainder <= num_q % den_q;
        end else begin
            quotient  <= {48{1'b1}};         // 除零: 全 1
            remainder <= num_q;
        end
        done <= 1'b1; busy <= 1'b0;
    end else wait_cnt <= wait_cnt + 1;
end
\end{Verilog}

两个设计点：\textbf{16 拍延迟}是为了给下游状态机"有等待可等"
（AWB 的控制环本来就按若干拍设计）；\textbf{除零给全 1}是约定的
哨兵值。再次强调：TD 综合时这个文件\textbf{不进工程}，
真网表在位。

\subsection{ph1p\_logic\_eram\_stub：双口 RAM 行为模型}

官方 blk\_mem 生成的 RTL 例化 TD 原语，参数/端口极多。
替身的编写流程是"从实例化现场抄"：

\begin{center}
\begin{tabular}{ll}
\toprule
\rowcolor{rowtint}\textbf{项目} & \textbf{处理} \\
\midrule
参数（DATA\_WIDTH\_A 等） & 全部声明，多数不影响行为 \\
端口 clka/clkb/cea/ceb & 时钟与使能，行为使用 \\
端口 dia/dib/doa/dob & 数据 \\
端口 bea/beb & 字节使能（全字写时忽略） \\
\textbf{ECC 端口} & sbiterr/dbiterr 输出；*inj 输入 \\
\bottomrule
\end{tabular}
\end{center}

\begin{warnbox}{ECC 端口方向抄反过一次}
第一版替身按"错误信号应该是输入吧"想当然，把 \code{sbiterr/dbiterr}
写成输入——ModelSim 编译通过（方向错误在 Verilog 端口连接里
不总是报错），但上板/综合语义全乱。修复：回官方实例化清单
逐个核对方向。\textbf{替身的价值 = 接口逐位忠实，不是"能编过"}。
\end{warnbox}

% ============================================================
\section{system\_top 逐段}

\subsection{信号流与"三个域"}

\begin{Verilog}
I_clk (axi4s_video_aclk, 74.25M 级):  ISP / 视觉 / 遥测 / 声光
I_pwm_clk (50MHz):                    gimbal_ctrl / pwm_servo
两者在 vision_top 内部经 pulse_sync/cdc_4phase 跨域(视觉已备)
\end{Verilog}

顶层本身\textbf{不做任何跨域}——跨域全部封装在 vision\_top 内部，
顶层只负责摆位置、接按键、引输出。这是"分层"的又一处体现。

\subsection{按键 glue 与遥测的层次化引用}

\begin{Verilog}
if (key0_trig) preset_idx <= preset_idx + 3'd1;   // KEY0: 阈值组
if (key1_trig) track_en   <= ~track_en;           // KEY1: 跟踪使能

uart_telemetry u_uart (
    .I_frame_tick (u_vision.t_tick),   // 层次化引用: 顶层特权
    .I_target_valid(u_vision.t_valid),
    .I_cx (u_vision.t_cx), .I_cy (u_vision.t_cy), ...);
\end{Verilog}

遥测拿的是视觉内部信号（\code{t\_tick/t\_valid/t\_cx/t\_cy}）。
选层次化引用而不是加端口：\textbf{遥测是"观测者"，不该改变被观测
对象的接口}——视觉模块作为交付给上板的产物保持干净；
观测通道只在集成顶层存在。

\subsection{仿真与上板的双用语法}

顶层文件同时服务两个场景：ModelSim 仿真与 TD 对接。
它\textbf{不含}三种黑盒（DDR2 / MIPI-PHY / HDMI）——仿真里
以帧转储代替 HDMI，DDR 写流空置；上板时这些角色由官方
wrapper 承担，本文件的视觉挂载段整体搬过去。

% ============================================================
\section{测试平台设计}

\subsection{虚拟传感器生成器}

\begin{Verilog}
raw = 512 灰底 (10bit)
红块 80x80 @ (500,300): R 位点=900, G/B 位点=64
噪声: 0.2% 像素替换为随机值 (default_rng(500) 固定种子)
打包: 40bit = {P0[39:30], P1[29:20], P2[19:10], P3[9:0]}
标志: 每行第 0 拍 sof、第 319 拍 eol, 单独文件一行一值
\end{Verilog}

红块中心 (540, 340) 与视觉 TPG 测试的红块位置接近——
整链仿真复用视觉侧已有的期望坐标口径（±40px）。

\subsection{TB 的四个检查器}

\begin{itemize}
  \item \textbf{几何检查器}：在 isp\_tvalid 上数 tuser/tlast。
    tuser 到达即"结账"上一帧（tuser 恰 1、tlast 恰 720），
    结账不符打 \code{GEO\_ERR}；
  \item \textbf{检测检查器}：\code{t\_upd} 拍采样锁定结果，
    连续多帧坐标在框内则 \code{det\_ok} 计数；
  \item \textbf{UART 解码器}：一个 9 态状态机按 A5..5A 组帧，
    校验和 + valid + 坐标范围三项对齐才计 \code{u\_good}；
  \item \textbf{终判与超时}：帧数够、检测够、UART 够则 PASS；
    450ms 超时 FAIL 并打印各计数——\textbf{失败也要给出诊断数据}。
\end{itemize}

\begin{keybox}{为什么 TB 里要有 UART 解码器}
外设层（遥测）是唯一能"整链端到端"观测检测结果的通道——
屏幕上看到的是 OSD，日志里看到的是 UART。两个通道互为佐证，
也提前验证了第二册的遥测协议。判据 D 用的是同一个
\code{uart\_rx\_mon}，一个监视器两处复用。
\end{keybox}

% ============================================================
\section{迭代复盘（本册相关 5 轮）}

\begin{center}
\begin{tabular}{llp{5.0cm}l}
\toprule
\rowcolor{rowtint}\textbf{轮} & \textbf{现象} & \textbf{根因 / 处置} & \textbf{类型} \\
\midrule
1 & elaboration 58840 错误 & 门级网表 + TD 原语 + 缺拷贝；
  写两个行为替身 + 补 data\_96bit\_to\_128bit & 环境 \\
2 & I\_tready 接常量报错 & 端口是输出，命名反向 & 集成 \\
3 & 重复声明 & 先赋值后声明，隐式 wire 冲突 & RTL \\
7 & 整链 tuser 全灭 & 官方 demosaic 依赖真实 CSI 时序；
  列 TD 阶段项 + 检查点文档化 & 界碑 \\
8 & RAW10 生成器只出 1 帧 & heredoc 转义吞替换；
  改行级补丁修复（4 帧 921600 拍） & 脚本 \\
\bottomrule
\end{tabular}
\end{center}

\begin{warnbox}{轮 1 的启示：58840 个错误不是"改 58840 次"}
三个根因、两个替身、一次补拷贝——\textbf{错误数量的规模不代表工作量}
。遇到海量报错先聚类：58840 条里其实只有 3 类。这个习惯在
RTL 调试里到处适用。
\end{warnbox}

\begin{tipbox}{脚本类 bug 的通用教训（轮 8）}
生成器出错时不要用"更复杂的转义"去绕——直接\textbf{改行级补丁}
（读行、替换指定行、写回）。本仓库为此吃过三次亏
（Step2 的 fix 脚本改坏 tex、本轮 heredoc 吞替换），
结论写进方法论：\textbf{改文件用行级操作，别用嵌套转义。}
\end{tipbox}

% ============================================================
\section{上板检查点与已知差异（交接给上板阶段）}

\begin{center}
\begin{tabular}{ll}
\toprule
\rowcolor{rowtint}\textbf{现象（上板若遇到）} & \textbf{先查这里} \\
\midrule
有数据无帧同步 & demosaic tuser 链（本仿真不激活的同一问题） \\
时序不收敛 & divider\_gate 门级网表路径；例程本身未收敛为已知现象 \\
跟踪颜色漂移 & AWB 自适应（对策：锁曝光/切阈值组 8 组） \\
画面随舵机动作闪 & 舵机供电（必须独立 5V 共地，仿真无法暴露） \\
\bottomrule
\end{tabular}
\end{center}

\begin{okbox}{本册交付清单}
\begin{itemize}
  \item \code{system\_top.v}：集成顶层（elaboration 零错误）
  \item 两个行为替身 + filelist 的"仿真专用"清单
  \item 虚拟 RAW10 传感器生成器（固定种子可复现）
  \item \code{tb\_system\_integration}：四判据 + 超时诊断
  \item 分层证明材料与上板检查点清单
\end{itemize}
\end{okbox}


% ============================================================
\part*{结语}
\clearpage
\section*{总结与下一步}
\addcontentsline{toc}{section}{总结与下一步}

\subsection*{本册学到了什么}
\addcontentsline{toc}{subsection}{本册学到了什么}
\begin{itemize}
  \item 官方 ISP 链路逐级职责（RAW10/BLC/demosaic/AWB）与明码/黑盒边界
  \item 96bit 契约三要素与 AWB 后挂载的三个理由
  \item 行为级替身的三条纪律（端口对照/行为对齐/禁止进综合）
  \item system\_top 集成与两处真实踩坑（端口方向、先赋值后声明）
  \item 虚拟 RAW10 传感器与整链四判据（功能级验收口径）
  \item 分层证明的写法与 demosaic tuser 界碑的处置
\end{itemize}

\subsection*{下一步学什么}
\addcontentsline{toc}{subsection}{下一步学什么}
\begin{itemize}
  \item \textbf{第二册：外设驱动、摄像头配置与固化}
  \begin{itemize}
    \item UART 遥测协议与逐字节验证
    \item I2C 总线、SC500CS 寄存器表与配置链
    \item 按键消抖与 TD 工程组织、QSPI 固化 runbook
  \end{itemize}
  \item 对应目录：\code{04\_embedded/03\_project/}
  \item 接口文档：\code{06\_interfaces/current/}（IF-MIPI / IF-DDR / IF-VISION）
\end{itemize}

\end{document}
